% --- Archon physics formalization source begin ---
% archon:physics
% archon:covers PhyXMiniProblems/problem_phyx_mini_0544.lean
% archon:source-report reports/phyx_mini/problem_phyx_mini_0544.source.json

\chapter{Physics problem phyx\_mini\_0544}
\label{ch:PhyXMiniProblems_problem_phyx_mini_0544}

\paragraph{Problem source.}
\#\# Physical scenario

Consider a potential energy step as shown in figure with a beam of particles incident from the left.

\#\# Question

Calculate the reflection coefficient for the case where the energy of the incident particles is greater than the height of the step

\#\# Answer choices

- A: A: \textbackslash{}(\textbackslash{}frac\{3 \textbackslash{}sqrt\{E\} \textbackslash{}sqrt\{E - V\_0\}\}\{(\textbackslash{}sqrt\{E\} + \textbackslash{}sqrt\{E - V\_0\})\textasciicircum{}2\}\textbackslash{})

- B: B: \textbackslash{}(\textbackslash{}frac\{4 \textbackslash{}sqrt\{E\} \textbackslash{}sqrt\{E - V\_0\}\}\{(\textbackslash{}sqrt\{E\} + \textbackslash{}sqrt\{E + V\_0\})\textasciicircum{}2\}\textbackslash{})

- C: C: \textbackslash{}(\textbackslash{}frac\{3 \textbackslash{}sqrt\{E\} \textbackslash{}sqrt\{E - V\_0\}\}\{(\textbackslash{}sqrt\{E\} + \textbackslash{}sqrt\{E + V\_0\})\textasciicircum{}2\}\textbackslash{})

- D: D: \textbackslash{}(\textbackslash{}frac\{4 \textbackslash{}sqrt\{E\} \textbackslash{}sqrt\{E - V\_0\}\}\{(\textbackslash{}sqrt\{E\} + \textbackslash{}sqrt\{E - V\_0\})\textasciicircum{}2\}\textbackslash{})

\#\# Figure caption (auxiliary; use the image as primary evidence)

The image depicts a potential energy diagram, typically used in quantum mechanics to represent a step potential. It features a graph with two axes. The horizontal axis is labeled \textbackslash{}(x\textbackslash{}), and the vertical axis is labeled \textbackslash{}(V(x)\textbackslash{}).

- From \textbackslash{}(x = -\textbackslash{}infty\textbackslash{}) to \textbackslash{}(x = 0\textbackslash{}), the potential \textbackslash{}(V(x)\textbackslash{}) is zero.
- At \textbackslash{}(x = 0\textbackslash{}), there is a sharp vertical increase in the potential.
- For \textbackslash{}(x > 0\textbackslash{}), the potential is constant at a higher level, labeled as \textbackslash{}(V\_0\textbackslash{}).

This represents a step potential where the potential energy sharply increases and remains constant beyond a certain point on the x-axis.

\#\# Dataset metadata

Quantum Phenomena; Physical Model Grounding Reasoning, Numerical Reasoning

\paragraph{Recorded answer/context.}
D: D: \textbackslash{}(\textbackslash{}frac\{4 \textbackslash{}sqrt\{E\} \textbackslash{}sqrt\{E - V\_0\}\}\{(\textbackslash{}sqrt\{E\} + \textbackslash{}sqrt\{E - V\_0\})\textasciicircum{}2\}\textbackslash{})

\paragraph{Figure/image path.}
/root/proposal\_for\_physic/science-mango/phyx\_mini\_run/phyx\_data/test\_image/544.png

\paragraph{Formalization target.}
create a compiling Lean file with sorry bodies at `PhyXMiniProblems/problem\_phyx\_mini\_0544.lean`. The Lean declarations must preserve the physical quantities, dimensions or dimensional roles, figure labels, governing-law hypotheses, and final relation expressed by this problem.
Use Mathlib/Physlib names found through LeanExplore where available. If a physics API is missing, introduce faithful local abstractions rather than scalar placeholder aliases.

\begin{theorem}[Physics formalization target]
\label{thm:physics:phyx_mini_0544:target}
\lean{PhyXMiniProblems.ProblemPhyXMini0544.problem_phyx_mini_0544}
\uses{def:physics:phyx-mini-0544:phyxminiproblems-problemphyxmini0544-potentialstepscatteringsetup, def:physics:phyx-mini-0544:phyxminiproblems-problemphyxmini0544-incidentenergyjoules, def:physics:phyx-mini-0544:phyxminiproblems-problemphyxmini0544-stepheightjoules, def:physics:phyx-mini-0544:phyxminiproblems-problemphyxmini0544-matchesproblemscenario, def:physics:phyx-mini-0544:phyxminiproblems-problemphyxmini0544-matchesprimarypotentialstepfigure, def:physics:phyx-mini-0544:phyxminiproblems-problemphyxmini0544-hasphysicalabovestepparameters, def:physics:phyx-mini-0544:phyxminiproblems-problemphyxmini0544-satisfiesabovestepdispersionlaw, def:physics:phyx-mini-0544:phyxminiproblems-problemphyxmini0544-satisfiesstepboundarymatchinglaws, def:physics:phyx-mini-0544:phyxminiproblems-problemphyxmini0544-matchesanswerchoice, def:physics:phyx-mini-0544:phyxminiproblems-problemphyxmini0544-recordeddatasetanswer}
The assigned autoformalize agent should translate this physics problem into Lean declarations in the covered file, with theorem and lemma proof bodies written as `by sorry`.
\end{theorem}
\begin{proof}
This is an autoformalization task, not a proof task. Produce statements that can later be proved by the physics prover without weakening the physical model.
\end{proof}

% --- Archon declaration topology begin ---
\section{Declaration topology}
\begin{definition}[Wave Number Quantity]
  \label{def:physics:phyx-mini-0544:phyxminiproblems-problemphyxmini0544-wavenumberquantity}
  \lean{PhyXMiniProblems.ProblemPhyXMini0544.WaveNumberQuantity}
  A unit-independent one-dimensional wave number, of dimension length⁻¹
\end{definition}
\begin{proof}
  This is the declared model definition of Wave Number Quantity.
\end{proof}
\begin{definition}[energy In Joules]
  \label{def:physics:phyx-mini-0544:phyxminiproblems-problemphyxmini0544-energyinjoules}
  \lean{PhyXMiniProblems.ProblemPhyXMini0544.energyInJoules}
  Coherent-SI readout of an energy, in joules
\end{definition}
\begin{proof}
  This is the declared model definition of energy In Joules.
\end{proof}
\begin{definition}[wave Number In Inverse Meters]
  \label{def:physics:phyx-mini-0544:phyxminiproblems-problemphyxmini0544-wavenumberininversemeters}
  \lean{PhyXMiniProblems.ProblemPhyXMini0544.waveNumberInInverseMeters}
  \uses{def:physics:phyx-mini-0544:phyxminiproblems-problemphyxmini0544-wavenumberquantity}
  Coherent-SI readout of a wave number, in inverse metres
\end{definition}
\begin{proof}
  This is the declared model definition of wave Number In Inverse Meters.
\end{proof}
\begin{definition}[Horizontal Direction]
  \label{def:physics:phyx-mini-0544:phyxminiproblems-problemphyxmini0544-horizontaldirection}
  \lean{PhyXMiniProblems.ProblemPhyXMini0544.HorizontalDirection}
  Horizontal propagation directions for the particle beam
\end{definition}
\begin{proof}
  This is the declared model definition of Horizontal Direction.
\end{proof}
\begin{definition}[Interface Side]
  \label{def:physics:phyx-mini-0544:phyxminiproblems-problemphyxmini0544-interfaceside}
  \lean{PhyXMiniProblems.ProblemPhyXMini0544.InterfaceSide}
  The side of the interface from which a beam approaches
\end{definition}
\begin{proof}
  This is the declared model definition of Interface Side.
\end{proof}
\begin{definition}[Incident Particle Beam]
  \label{def:physics:phyx-mini-0544:phyxminiproblems-problemphyxmini0544-incidentparticlebeam}
  \lean{PhyXMiniProblems.ProblemPhyXMini0544.IncidentParticleBeam}
  \uses{def:physics:phyx-mini-0544:phyxminiproblems-problemphyxmini0544-horizontaldirection, def:physics:phyx-mini-0544:phyxminiproblems-problemphyxmini0544-interfaceside}
  The particle beam described in the prose
\end{definition}
\begin{proof}
  This is the declared model definition of Incident Particle Beam.
\end{proof}
\begin{definition}[Figure Axis]
  \label{def:physics:phyx-mini-0544:phyxminiproblems-problemphyxmini0544-figureaxis}
  \lean{PhyXMiniProblems.ProblemPhyXMini0544.FigureAxis}
  Horizontal and vertical axes in the supplied potential-energy diagram
\end{definition}
\begin{proof}
  This is the declared model definition of Figure Axis.
\end{proof}
\begin{definition}[Axis Symbol]
  \label{def:physics:phyx-mini-0544:phyxminiproblems-problemphyxmini0544-axissymbol}
  \lean{PhyXMiniProblems.ProblemPhyXMini0544.AxisSymbol}
  Literal symbols printed at the ends of the two axes
\end{definition}
\begin{proof}
  This is the declared model definition of Axis Symbol.
\end{proof}
\begin{definition}[Potential Level Label]
  \label{def:physics:phyx-mini-0544:phyxminiproblems-problemphyxmini0544-potentiallevellabel}
  \lean{PhyXMiniProblems.ProblemPhyXMini0544.PotentialLevelLabel}
  Potential-level symbols printed in the image
\end{definition}
\begin{proof}
  This is the declared model definition of Potential Level Label.
\end{proof}
\begin{definition}[Potential Step Figure]
  \label{def:physics:phyx-mini-0544:phyxminiproblems-problemphyxmini0544-potentialstepfigure}
  \lean{PhyXMiniProblems.ProblemPhyXMini0544.PotentialStepFigure}
  \uses{def:physics:phyx-mini-0544:phyxminiproblems-problemphyxmini0544-figureaxis, def:physics:phyx-mini-0544:phyxminiproblems-problemphyxmini0544-axissymbol, def:physics:phyx-mini-0544:phyxminiproblems-problemphyxmini0544-potentiallevellabel}
  Raster-visible data from image 544. Booleans record the qualitative line segments; the interface coordinate is the displayed scalar label 0
\end{definition}
\begin{proof}
  This is the declared model definition of Potential Step Figure.
\end{proof}
\begin{definition}[Potential Step Scattering Setup]
  \label{def:physics:phyx-mini-0544:phyxminiproblems-problemphyxmini0544-potentialstepscatteringsetup}
  \lean{PhyXMiniProblems.ProblemPhyXMini0544.PotentialStepScatteringSetup}
  \uses{def:physics:phyx-mini-0544:phyxminiproblems-problemphyxmini0544-wavenumberquantity, def:physics:phyx-mini-0544:phyxminiproblems-problemphyxmini0544-incidentparticlebeam, def:physics:phyx-mini-0544:phyxminiproblems-problemphyxmini0544-potentialstepfigure}
  The scattering setup. The potential and step height are physical energies. The three real amplitudes are the normalized coefficients of the incident, reflected, and transmitted stationary plane-wave components at the interface
\end{definition}
\begin{proof}
  This is the declared model definition of Potential Step Scattering Setup.
\end{proof}
\begin{definition}[incident Energy Joules]
  \label{def:physics:phyx-mini-0544:phyxminiproblems-problemphyxmini0544-incidentenergyjoules}
  \lean{PhyXMiniProblems.ProblemPhyXMini0544.incidentEnergyJoules}
  \uses{def:physics:phyx-mini-0544:phyxminiproblems-problemphyxmini0544-energyinjoules, def:physics:phyx-mini-0544:phyxminiproblems-problemphyxmini0544-potentialstepscatteringsetup}
  Joule readout of the incident particle energy E
\end{definition}
\begin{proof}
  This is the declared model definition of incident Energy Joules.
\end{proof}
\begin{definition}[step Height Joules]
  \label{def:physics:phyx-mini-0544:phyxminiproblems-problemphyxmini0544-stepheightjoules}
  \lean{PhyXMiniProblems.ProblemPhyXMini0544.stepHeightJoules}
  \uses{def:physics:phyx-mini-0544:phyxminiproblems-problemphyxmini0544-energyinjoules, def:physics:phyx-mini-0544:phyxminiproblems-problemphyxmini0544-potentialstepscatteringsetup}
  Joule readout of the step height V₀
\end{definition}
\begin{proof}
  This is the declared model definition of step Height Joules.
\end{proof}
\begin{definition}[incident Wave Number SI]
  \label{def:physics:phyx-mini-0544:phyxminiproblems-problemphyxmini0544-incidentwavenumbersi}
  \lean{PhyXMiniProblems.ProblemPhyXMini0544.incidentWaveNumberSI}
  \uses{def:physics:phyx-mini-0544:phyxminiproblems-problemphyxmini0544-wavenumberininversemeters, def:physics:phyx-mini-0544:phyxminiproblems-problemphyxmini0544-potentialstepscatteringsetup}
  Left-region wave number k₁, in inverse metres
\end{definition}
\begin{proof}
  This is the declared model definition of incident Wave Number SI.
\end{proof}
\begin{definition}[transmitted Wave Number SI]
  \label{def:physics:phyx-mini-0544:phyxminiproblems-problemphyxmini0544-transmittedwavenumbersi}
  \lean{PhyXMiniProblems.ProblemPhyXMini0544.transmittedWaveNumberSI}
  \uses{def:physics:phyx-mini-0544:phyxminiproblems-problemphyxmini0544-wavenumberininversemeters, def:physics:phyx-mini-0544:phyxminiproblems-problemphyxmini0544-potentialstepscatteringsetup}
  Right-region wave number k₂, in inverse metres
\end{definition}
\begin{proof}
  This is the declared model definition of transmitted Wave Number SI.
\end{proof}
\begin{definition}[Matches Problem Scenario]
  \label{def:physics:phyx-mini-0544:phyxminiproblems-problemphyxmini0544-matchesproblemscenario}
  \lean{PhyXMiniProblems.ProblemPhyXMini0544.MatchesProblemScenario}
  \uses{def:physics:phyx-mini-0544:phyxminiproblems-problemphyxmini0544-potentialstepscatteringsetup}
  The incident beam approaches the step from the left and travels rightward
\end{definition}
\begin{proof}
  This is the declared model definition of Matches Problem Scenario.
\end{proof}
\begin{definition}[Matches Primary Potential Step Figure]
  \label{def:physics:phyx-mini-0544:phyxminiproblems-problemphyxmini0544-matchesprimarypotentialstepfigure}
  \lean{PhyXMiniProblems.ProblemPhyXMini0544.MatchesPrimaryPotentialStepFigure}
  \uses{def:physics:phyx-mini-0544:phyxminiproblems-problemphyxmini0544-energyinjoules, def:physics:phyx-mini-0544:phyxminiproblems-problemphyxmini0544-potentialstepscatteringsetup, def:physics:phyx-mini-0544:phyxminiproblems-problemphyxmini0544-stepheightjoules}
  All geometric and potential-energy readouts supplied by the primary bitmap. The value exactly at the discontinuity is intentionally left unspecified
\end{definition}
\begin{proof}
  This is the declared model definition of Matches Primary Potential Step Figure.
\end{proof}
\begin{definition}[Has Physical Above Step Parameters]
  \label{def:physics:phyx-mini-0544:phyxminiproblems-problemphyxmini0544-hasphysicalabovestepparameters}
  \lean{PhyXMiniProblems.ProblemPhyXMini0544.HasPhysicalAboveStepParameters}
  \uses{def:physics:phyx-mini-0544:phyxminiproblems-problemphyxmini0544-potentialstepscatteringsetup, def:physics:phyx-mini-0544:phyxminiproblems-problemphyxmini0544-incidentenergyjoules, def:physics:phyx-mini-0544:phyxminiproblems-problemphyxmini0544-stepheightjoules, def:physics:phyx-mini-0544:phyxminiproblems-problemphyxmini0544-incidentwavenumbersi, def:physics:phyx-mini-0544:phyxminiproblems-problemphyxmini0544-transmittedwavenumbersi}
  The positive, propagating branch specified by E > V₀
\end{definition}
\begin{proof}
  This is the declared model definition of Has Physical Above Step Parameters.
\end{proof}
\begin{definition}[Satisfies Above Step Dispersion Law]
  \label{def:physics:phyx-mini-0544:phyxminiproblems-problemphyxmini0544-satisfiesabovestepdispersionlaw}
  \lean{PhyXMiniProblems.ProblemPhyXMini0544.SatisfiesAboveStepDispersionLaw}
  \uses{def:physics:phyx-mini-0544:phyxminiproblems-problemphyxmini0544-potentialstepscatteringsetup, def:physics:phyx-mini-0544:phyxminiproblems-problemphyxmini0544-incidentenergyjoules, def:physics:phyx-mini-0544:phyxminiproblems-problemphyxmini0544-stepheightjoules, def:physics:phyx-mini-0544:phyxminiproblems-problemphyxmini0544-incidentwavenumbersi, def:physics:phyx-mini-0544:phyxminiproblems-problemphyxmini0544-transmittedwavenumbersi}
  For a fixed particle mass and reduced Planck constant, the nonrelativistic wave number is a common positive multiple of the square root of the kinetic energy. The existential scale has coherent-SI units m⁻¹ J⁻¹/²; introducing it avoids representing mass or ℏ, neither of which is a stated problem parameter
\end{definition}
\begin{proof}
  This is the declared model definition of Satisfies Above Step Dispersion Law.
\end{proof}
\begin{definition}[Satisfies Step Boundary Matching Laws]
  \label{def:physics:phyx-mini-0544:phyxminiproblems-problemphyxmini0544-satisfiesstepboundarymatchinglaws}
  \lean{PhyXMiniProblems.ProblemPhyXMini0544.SatisfiesStepBoundaryMatchingLaws}
  \uses{def:physics:phyx-mini-0544:phyxminiproblems-problemphyxmini0544-potentialstepscatteringsetup, def:physics:phyx-mini-0544:phyxminiproblems-problemphyxmini0544-incidentwavenumbersi, def:physics:phyx-mini-0544:phyxminiproblems-problemphyxmini0544-transmittedwavenumbersi}
  Continuity of a stationary state and its derivative at a finite potential step gives the first two fields. The final field states the flux-ratio interpretation used by the recorded expression: right-moving transmitted flux divided by incident flux. This is a governing relation in amplitudes and wave numbers, not the requested energy-only closed form
\end{definition}
\begin{proof}
  This is the declared model definition of Satisfies Step Boundary Matching Laws.
\end{proof}
\begin{definition}[Answer Choice]
  \label{def:physics:phyx-mini-0544:phyxminiproblems-problemphyxmini0544-answerchoice}
  \lean{PhyXMiniProblems.ProblemPhyXMini0544.AnswerChoice}
  Labels of the four multiple-choice expressions in the source
\end{definition}
\begin{proof}
  This is the declared model definition of Answer Choice.
\end{proof}
\begin{definition}[displayed Coefficient]
  \label{def:physics:phyx-mini-0544:phyxminiproblems-problemphyxmini0544-displayedcoefficient}
  \lean{PhyXMiniProblems.ProblemPhyXMini0544.displayedCoefficient}
  \uses{def:physics:phyx-mini-0544:phyxminiproblems-problemphyxmini0544-answerchoice}
  The coefficient expression printed beside each answer label
\end{definition}
\begin{proof}
  This is the declared model definition of displayed Coefficient.
\end{proof}
\begin{definition}[Matches Answer Choice]
  \label{def:physics:phyx-mini-0544:phyxminiproblems-problemphyxmini0544-matchesanswerchoice}
  \lean{PhyXMiniProblems.ProblemPhyXMini0544.MatchesAnswerChoice}
  \uses{def:physics:phyx-mini-0544:phyxminiproblems-problemphyxmini0544-potentialstepscatteringsetup, def:physics:phyx-mini-0544:phyxminiproblems-problemphyxmini0544-incidentenergyjoules, def:physics:phyx-mini-0544:phyxminiproblems-problemphyxmini0544-stepheightjoules, def:physics:phyx-mini-0544:phyxminiproblems-problemphyxmini0544-answerchoice, def:physics:phyx-mini-0544:phyxminiproblems-problemphyxmini0544-displayedcoefficient}
  The modeled coefficient agrees with the expression at a displayed label
\end{definition}
\begin{proof}
  This is the declared model definition of Matches Answer Choice.
\end{proof}
\begin{definition}[recorded Dataset Answer]
  \label{def:physics:phyx-mini-0544:phyxminiproblems-problemphyxmini0544-recordeddatasetanswer}
  \lean{PhyXMiniProblems.ProblemPhyXMini0544.recordedDatasetAnswer}
  \uses{def:physics:phyx-mini-0544:phyxminiproblems-problemphyxmini0544-answerchoice}
  Answer label recorded by the dataset; this is source metadata, not a premise
\end{definition}
\begin{proof}
  This is the declared model definition of recorded Dataset Answer.
\end{proof}
\begin{lemma}[reflection Coefficient energy formula]
  \label{lem:physics:phyx-mini-0544:phyxminiproblems-problemphyxmini0544-reflectioncoefficient-energy-formula}
  \lean{PhyXMiniProblems.ProblemPhyXMini0544.reflectionCoefficient_energy_formula}
  \uses{def:physics:phyx-mini-0544:phyxminiproblems-problemphyxmini0544-potentialstepscatteringsetup, def:physics:phyx-mini-0544:phyxminiproblems-problemphyxmini0544-incidentenergyjoules, def:physics:phyx-mini-0544:phyxminiproblems-problemphyxmini0544-stepheightjoules, def:physics:phyx-mini-0544:phyxminiproblems-problemphyxmini0544-hasphysicalabovestepparameters, def:physics:phyx-mini-0544:phyxminiproblems-problemphyxmini0544-satisfiesabovestepdispersionlaw, def:physics:phyx-mini-0544:phyxminiproblems-problemphyxmini0544-satisfiesstepboundarymatchinglaws}
  The dispersion and boundary-matching laws yield the source's energy-only coefficient expression. This derived statement is not included in any premise structure
\end{lemma}
\begin{proof}
  The conclusion follows from the governing relations and figure readouts named in the statement of reflection Coefficient energy formula.
\end{proof}
% --- Archon declaration topology end ---
% --- Archon physics formalization source end ---
